\chapter{Probabilistic Foundations and Limit Theorems}

\section{Probability and Moment Bounds}

\subsection{Foundations of Probability Theory}

To discuss limit theorems rigorously, we must first establish the measure-theoretic foundation of probability. This framework prevents the paradoxes that arise in naive probability and allows us to handle continuous and infinite-dimensional spaces seamlessly.

\begin{definition}[Probability Space]
A \textit{probability space} is a triple $(\Omega, \mathcal{F}, \Prob)$ consisting of:
\begin{enumerate}
\item \textit{Sample space ($\Omega$)}: a nonempty set of all possible outcomes of a random experiment. Each element $\omega$ represents one individual outcome.
\item \textit{$\sigma$-algebra of events ($\mathcal{F}$)}: a collection of subsets of $\Omega$, called events, that contains $\Omega$ and is closed under complements relative to $\Omega$ and countable unions. Explicitly, for every event $A$ and every sequence of events $A_n$, these conditions are
\[
\begin{gathered}
\Omega \in \mathcal{F}, \\
\Omega \setminus A \in \mathcal{F}, \\
\bigcup_{n=1}^{\infty} A_n \in \mathcal{F}.
\end{gathered}
\]
\item \textit{Probability measure ($\Prob$)}: a function assigning to each event a number between zero and one, with total probability one:
\[
\begin{gathered}
\Prob: \mathcal{F} \to [0,1], \\
\Prob(\Omega) = 1.
\end{gathered}
\]
It is countably additive: for every sequence of pairwise disjoint events $A_n$, the probability of their union equals the sum of their probabilities:
\[
\Prob\!\left(\bigcup_{n=1}^{\infty} A_n\right)
= \sum_{n=1}^{\infty} \Prob(A_n).
\]
\end{enumerate}
\end{definition}

\begin{remark}[Why a $\sigma$-algebra?]
We cannot simply assign probabilities to \textit{all} subsets of $\Omega$ (for example, if $\Omega = \R$, the Axiom of Choice implies the existence of non-measurable sets like the Vitali set). The $\sigma$-algebra $\mathcal{F}$ restricts the domain of $\Prob$ to a collection of ``well-behaved'' sets. Closure under countable unions is strictly necessary because probability theory frequently involves limits of events (e.g., ``the event happens eventually''), and we must guarantee that the limit of a sequence of valid events is still a valid event to which we can assign a probability.
\end{remark}

\begin{definition}[Random Variable and Expectation]
A \textit{random variable} $S$ is a measurable function
\[
S: (\Omega, \mathcal{F}) \to (\R, \mathcal{B}(\R)),
\]
where $\mathcal{B}(\R)$ is the Borel $\sigma$-algebra on $\R$. Measurability means that for any Borel set $B$, the preimage
\[
S^{-1}(B) = \{\omega \in \Omega : S(\omega) \in B\}
\]
is an element of $\mathcal{F}$.
For the purposes of this chapter, we restrict attention to integrable
random variables, meaning that
\[
\E[|S|]=\int_\Omega |S(\omega)|\,\dd\Prob(\omega)<\infty.
\]
For such a random variable, the \textit{expectation} (or expected value)
is the Lebesgue integral
\begin{equation}
\E[S] = \int_\Omega S(\omega) \dd \Prob(\omega)
\end{equation}
and is finite. More generally, extended expectations can be defined from
the positive and negative parts when they are not both infinite.
\end{definition}

\begin{remark}[Why Measurability?]
We require $S$ to be measurable so that we can calculate probabilities of events defined by $S$. For instance, to compute $\Prob(S \leq x)$, we need the set $\{\omega : S(\omega) \leq x\}$ to be in $\mathcal{F}$ so that the measure $\Prob$ can be applied to it. Without measurability, the expectation integral would be undefined.
\end{remark}

Figure~\ref{fig:atlas-conditional-events} visualizes conditioning as restricting attention to a specified event.

\begin{figure}[!htbp]
\centering
\begin{tikzpicture}[scale=1.45,>=Stealth]
\draw[black!60] (-2.5,-1.6) rectangle (2.5,1.6);
\fill[figblue!15] (0.65,0) circle (1.25);
\begin{scope}
\clip (-0.65,0) circle (1.25);
\fill[figgreen!45] (0.65,0) circle (1.25);
\end{scope}
\draw[figred,thick] (-0.65,0) circle (1.25);
\draw[figblue,thick] (0.65,0) circle (1.25);
\node at (-1.25,0.25) {$A$};
\node at (1.25,0.25) {$B$};
\node at (0,0) {$A\cap B$};
\node at (2.2,1.3) {$\Omega$};
\end{tikzpicture}
\caption{The outer rectangle represents the sample space. Conditioning on $B$ restricts attention to the blue disk; its overlap with $A$ is highlighted in green. For $\Prob(B)>0$, the conditional probability is $\Prob(A\mid B)=\Prob(A\cap B)/\Prob(B)$. Areas are schematic and need not be proportional to probabilities.}
\label{fig:atlas-conditional-events}
\end{figure}

\subsection{Markov and Chebyshev Inequalities}

These inequalities provide deterministic bounds on the tails of probability distributions using only the moments (expectation and variance) of the random variable.

\begin{proposition}[Markov's Inequality]
Let $S$ be a nonnegative random variable ($S \geq 0$ almost surely) and let $a > 0$. Then:
\begin{equation}
\Prob(S \geq a) \leq \frac{\E[S]}{a} % <--- QUI USA \Prob
\end{equation}

\end{proposition}

\begin{proof}
We introduce the \textit{indicator function} $\mathbf{1}_{S \geq a}(\omega)$, which equals $1$ if $S(\omega) \geq a$ and $0$ otherwise.
Observe the pointwise inequality for all $\omega \in \Omega$:
\begin{equation}
S(\omega) \geq a \cdot \mathbf{1}_{S \geq a}(\omega)
\end{equation}

This holds because if $S(\omega) \geq a$, the right side is $a$, and $S(\omega) \geq a$; if $S(\omega) < a$, the right side is $0$, and $S(\omega) \geq 0$ by assumption.
Since the Lebesgue integral preserves inequalities for nonnegative functions, we integrate both sides with respect to $\Prob$ over $\Omega$:
\begin{equation}
\int_\Omega S(\omega) \dd \Prob(\omega) \geq \int_\Omega a \cdot \mathbf{1}_{S \geq a}(\omega) \dd \Prob(\omega)
\end{equation}

The left side is exactly $\E[S]$. On the right side, $a$ is a constant and can be pulled out. The integral of an indicator function over $\Omega$ is precisely the probability of the set it indicates:
\begin{equation}
\begin{aligned}
\E[S] &\geq a \int_\Omega \mathbf{1}_{S \geq a}(\omega) \dd \Prob(\omega) \\
&= a \cdot \Prob(S \geq a)
\end{aligned}
\end{equation}
Dividing both sides by $a > 0$ yields the result.
\end{proof}

\begin{definition}[Variance]
The \textit{variance} of a random variable $S$ with finite expectation $\mu = \E[S]$ is defined as the expected squared deviation from the mean:
\begin{equation}
\var(S) = \E[(S - \E[S])^2]
\end{equation}

\end{definition}

\begin{theorem}[Chebyshev's Inequality]
If $S$ has finite expectation $\mu$ and finite variance $\sigma^2 = \var(S) < \infty$, then for any $a > 0$:
\begin{equation}
\Prob(|S - \mu| \geq a) \leq \frac{\sigma^2}{a^2}
\end{equation}

\end{theorem}

\begin{proof}
Define a new random variable $X = (S - \mu)^2$. Since it is a square, $X \geq 0$ almost surely.
Notice that the event $\{|S - \mu| \geq a\}$ is mathematically equivalent to the event $\{(S - \mu)^2 \geq a^2\}$, which is exactly $\{X \geq a^2\}$.
We apply Markov's Inequality to the nonnegative random variable $X$ with the threshold $a^2 > 0$:
\begin{equation}
\Prob(X \geq a^2) \leq \frac{\E[X]}{a^2}
\end{equation}
Substituting back $X = (S - \mu)^2$ and recognizing that
\[
\begin{aligned}
\E[X] &= \E[(S - \mu)^2] \\
&= \var(S) \\
&= \sigma^2,
\end{aligned}
\]
we obtain:
\begin{equation}
\Prob(|S - \mu| \geq a) \leq \frac{\sigma^2}{a^2}
\end{equation}

\end{proof}

\section{Limit Theorems for Sums of Random Variables}

\subsection{Weak Law of Large Numbers}

The Weak Law of Large Numbers (WLLN) formalizes the intuitive idea that the average of a large number of independent observations converges to the expected value.

\begin{definition}[Independent and Identically Distributed (i.i.d.)]
A sequence of random variables $S_1, S_2, \dots$ is:
\begin{itemize}
\item \textit{Identically distributed}: if they all share the same cumulative distribution function, and consequently, the same expectation $\mu$ and variance $\sigma^2$.
\item \textit{Independent}: if the joint distribution of any finite subset is the product of their marginal distributions. Crucially for expectations, this implies that for $i \neq j$, the expectation of the product factors:

\[
\E[S_i S_j] = \E[S_i]\E[S_j].
\]
Consequently, their covariance is zero:
\[
\E[(S_i - \mu)(S_j - \mu)] = 0.
\]

\end{itemize}
\end{definition}

\begin{theorem}[Weak Law of Large Numbers]
Let $S_1, S_2, \dots$ be a sequence of i.i.d. random variables with mean $\mu$ and finite variance $\sigma^2$. Let the sample mean be
\[
z_m = (1/m) \sum_{i=1}^m S_i.
\]
Then for any $\varepsilon > 0$:
\begin{equation}
\lim_{m \to \infty} \Prob(|z_m - \mu| \geq \varepsilon) = 0
\end{equation}

\end{theorem}

\begin{proof}
We first explicitly calculate the expectation and variance of $z_m$.
By the linearity of the expectation operator:
\begin{equation}
\begin{aligned}
\E[z_m] &= \E\left[ \frac{1}{m} \sum_{i=1}^m S_i \right] \\
&= \frac{1}{m} \sum_{i=1}^m \E[S_i] \\
&= \frac{1}{m} (m \mu) \\
&= \mu
\end{aligned}
\end{equation}
Next, we calculate the variance of $z_m$. By definition:
\begin{equation}
\begin{aligned}
\var(z_m) &= \E[(z_m - \E[z_m])^2] \\
&= \E\left[ \left( \frac{1}{m} \sum_{i=1}^m (S_i - \mu) \right)^2 \right]
\end{aligned}
\end{equation}
Expanding the square of the sum yields a double sum:
\begin{equation}
\var(z_m) = \frac{1}{m^2} \E\left[ \sum_{i=1}^m \sum_{j=1}^m (S_i - \mu)(S_j - \mu) \right]
\end{equation}
Using the linearity of expectation again, we bring the expectation inside the sums:
\begin{equation}
\var(z_m) = \frac{1}{m^2} \sum_{i=1}^m \sum_{j=1}^m \E[(S_i - \mu)(S_j - \mu)]
\end{equation}
We evaluate the terms in the double sum based on the indices:
\begin{itemize}
\item \textit{Diagonal terms}: if $i = j$, the term is

\[
\begin{aligned}
\E[(S_i - \mu)^2] &= \var(S_i) \\
&= \sigma^2.
\end{aligned}
\]
There are exactly $m$ such terms on the diagonal of the double sum.
\item \textit{Off-diagonal terms}: if $i \neq j$, the term is $\E[(S_i - \mu)(S_j - \mu)]$. Because $S_i$ and $S_j$ are independent, this factors into

\[
\begin{aligned}
\E[S_i - \mu]\E[S_j - \mu] &= 0 \cdot 0 \\
&= 0.
\end{aligned}
\]
All $m^2 - m$ off-diagonal terms vanish.
\end{itemize}
Summing these up, we get:
\begin{equation}
\begin{aligned}
\var(z_m) &= \frac{1}{m^2} (m \sigma^2) \\
&= \frac{\sigma^2}{m}
\end{aligned}
\end{equation}
Now, we apply Chebyshev's Inequality to the random variable $z_m$ with threshold $\varepsilon > 0$:
\begin{equation}
\begin{aligned}
\Prob(|z_m - \mu| \geq \varepsilon) &\leq \frac{\var(z_m)}{\varepsilon^2} \\
&= \frac{\sigma^2}{m \varepsilon^2}
\end{aligned}
\end{equation}
Since $\sigma^2$ and $\varepsilon^2$ are fixed finite constants, taking the limit as $m \to \infty$ gives:
\begin{equation}
\lim_{m \to \infty} \frac{\sigma^2}{m \varepsilon^2} = 0
\end{equation}
Because probabilities are nonnegative, we conclude that
\[
\lim_{m \to \infty} \Prob(|z_m - \mu| \geq \varepsilon) = 0.
\]

\end{proof}

\subsection{Characteristic Functions and the Central Limit Theorem}


While the WLLN tells us the sample mean converges to a constant, the Central Limit Theorem (CLT) describes the exact distribution of the fluctuations around that constant.

\begin{definition}[Convergence in Distribution]
A sequence of random variables $Y_m$ \textit{converges in distribution} to a random variable $Y$, denoted $Y_m \xrightarrow{d} Y$, if their cumulative distribution functions (CDFs) $F_{Y_m}(y)$ converge to $F_Y(y)$ at every point $y$ where $F_Y(y)$ is continuous.
\end{definition}

\begin{remark}[Why use Characteristic Functions?]
Proving convergence of CDFs directly is analytically intractable. Instead, we use the \textit{Characteristic Function} (CF), which is the Fourier transform of the probability measure. The CF always exists, uniquely determines the probability distribution, and converts the convolution of independent random variables into simple algebraic multiplication, making limit calculations straightforward.
\end{remark}

\begin{definition}[Characteristic Function]
The characteristic function of a random variable $S$ is defined as:
\begin{equation}
\varphi_S(t) = \E[e^{itS}] = \int_\Omega e^{itS(\omega)} \dd \Prob(\omega), \quad t \in \R
\end{equation}

\end{definition}

\begin{theorem}[Lindeberg--Lévy Central Limit Theorem]
Let $S_1, S_2, \dots$ be a sequence of i.i.d. random variables with mean $\mu$ and finite variance $\sigma^2 > 0$. Let
\[
z_m = (1/m) \sum_{i=1}^m S_i.
\]

Then the normalized sample mean converges in distribution to a standard normal random variable:
\begin{equation}
Y_m = \frac{\sqrt{m}(z_m - \mu)}{\sigma} \xrightarrow{d} \mathcal{N}(0, 1)
\end{equation}

\end{theorem}

\subsection{Proof of the Central Limit Theorem}

\begin{proof}
We will prove this by showing that the characteristic function of $Y_m$ converges pointwise to the characteristic function of the standard normal distribution $\mathcal{N}(0,1)$.

\textit{Step 1: Express $Y_m$ as a sum of centered variables.}
Let $X_i = S_i - \mu$. The variables $X_i$ are i.i.d. with $\E[X_i] = 0$ and $\var(X_i) = \sigma^2$. We can rewrite $Y_m$ as:
\begin{equation}
\begin{aligned}
Y_m &= \frac{\sqrt{m}}{\sigma} \left( \frac{1}{m} \sum_{i=1}^m X_i \right) \\
&= \frac{1}{\sigma\sqrt{m}} \sum_{i=1}^m X_i
\end{aligned}
\end{equation}

\textit{Step 2: Relate the CF of $Y_m$ to the CF of $X_1$.}
Using the properties of characteristic functions:
\begin{itemize}
\item \textit{Scaling}: the following identity holds:

\[
\begin{aligned}
\varphi_{cX}(t) &= \E[e^{it(cX)}] \\
&= \varphi_X(ct).
\end{aligned}
\]

\item \textit{Independence}: if $X$ and $Y$ are independent,

\[
\begin{aligned}
\varphi_{X+Y}(t) &= \E[e^{itX}e^{itY}] \\
&= \E[e^{itX}]\E[e^{itY}] \\
&= \varphi_X(t)\varphi_Y(t).
\end{aligned}
\]

\end{itemize}
Applying these properties to $Y_m$:
\begin{equation}
\begin{aligned}
\varphi_{Y_m}(t) &= \varphi_{\sum X_i / (\sigma\sqrt{m})}(t) \\
&= \left[ \varphi_{X_1}\left( \frac{t}{\sigma\sqrt{m}} \right) \right]^m
\end{aligned}
\end{equation}

\textit{Step 3: Taylor expand the characteristic function of $X_1$.}
We expand $\varphi_{X_1}(u)$ around $u=0$. To do this rigorously, we differentiate under the integral sign. This is justified by the Dominated Convergence Theorem because the derivatives of the integrand are bounded by $|X_1|^k$, and we assume $\E[|X_1|^2] < \infty$.
\begin{align*}
\varphi_{X_1}(u) &= \E[e^{iuX_1}] \\
    \varphi_{X_1}'(u) &= \E[iX_1 e^{iuX_1}] \implies \varphi_{X_1}'(0) = i\E[X_1] = 0 \\
    \varphi_{X_1}''(u) &= \E[(iX_1)^2 e^{iuX_1}] = -\E[X_1^2 e^{iuX_1}] \implies \varphi_{X_1}''(0) = -\E[X_1^2] = -\var(X_1) = -\sigma^2
\end{align*}
By Taylor's theorem with the Peano form of the remainder, we have:
\begin{equation}
\begin{aligned}
\varphi_{X_1}(u) &= \varphi_{X_1}(0) + u\varphi_{X_1}'(0) + \frac{u^2}{2}\varphi_{X_1}''(0) + o(u^2) \\
&= 1 - \frac{\sigma^2 u^2}{2} + o(u^2)
\end{aligned}
\end{equation}
where $o(u^2)$ denotes a function such that
\[
\lim_{u \to 0} o(u^2)/u^2 = 0.
\]

\textit{Step 4: Substitute and take the limit.}
We substitute
\[
u = (t/(\sigma\sqrt{m}))
\]
into the Taylor expansion. Note that as $m \to \infty$, $u \to 0$:
\begin{equation}
\begin{aligned}
\varphi_{X_1}\left( \frac{t}{\sigma\sqrt{m}} \right) &= 1 - \frac{\sigma^2}{2} \left( \frac{t}{\sigma\sqrt{m}} \right)^2 + o\left( \frac{t^2}{\sigma^2 m} \right) \\
&= 1 - \frac{t^2}{2m} + o\left( \frac{1}{m} \right)
\end{aligned}
\end{equation}
Now, we substitute this back into the expression for $\varphi_{Y_m}(t)$:
\begin{equation}
\varphi_{Y_m}(t) = \left[ 1 - \frac{t^2}{2m} + o\left( \frac{1}{m} \right) \right]^m
\end{equation}

To evaluate the limit as $m \to \infty$, we rewrite the expression using the exponential and natural logarithm:
\begin{equation}
\varphi_{Y_m}(t) = \exp\left( m \ln\left( 1 - \frac{t^2}{2m} + o\left( \frac{1}{m} \right) \right) \right)
\end{equation}
Using the Taylor expansion $\ln(1 + x) = x + o(x)$ for small $x$, where
\[
x = -(t^2/(2m)) + o( (1/m) ):
\]

\begin{equation}
\begin{aligned}
m \ln\left( 1 - \frac{t^2}{2m} + o\left( \frac{1}{m} \right) \right) &= m \left( -\frac{t^2}{2m} + o\left( \frac{1}{m} \right) \right) \\
&= -\frac{t^2}{2} + m \cdot o\left( \frac{1}{m} \right)
\end{aligned}
\end{equation}
The remainder term satisfies
\[
\lim_{m \to \infty} m \cdot o(1/m) = 0,
\]
Therefore, the exponent converges to $-t^2/2$. Therefore:
\begin{equation}
\lim_{m \to \infty} \varphi_{Y_m}(t) = e^{-t^2/2}
\end{equation}

\textit{Step 5: Apply Lévy's Continuity Theorem.}
\textit{Lévy's Continuity Theorem} states that if a sequence of characteristic functions $\varphi_{Y_m}(t)$ converges pointwise to a limit function $\varphi(t)$ that is continuous at $t=0$, then $Y_m$ converges in distribution to a random variable $Y$ whose characteristic function is $\varphi(t)$.
The limit function $e^{-t^2/2}$ is continuous everywhere. Furthermore, it is exactly the characteristic function of the standard normal distribution $\mathcal{N}(0,1)$.
Thus, we conclude that $Y_m \xrightarrow{d} \mathcal{N}(0, 1)$. \end{proof}

\subsection{Normal Approximation and Interpretation}

Figure~\ref{fig:ch16-clt} illustrates the approach to a normal distribution described by the Central Limit Theorem:
\[
z_m \xrightarrow{d} N(0,1).
\]

\begin{figure}[!htbp]
\centering
\begin{tikzpicture}
\begin{axis}[
  width=13cm,height=8cm, axis lines=left,
  xlabel={$z=((\sqrt{m}(\bar X_m-\mu))/\sigma)$}, ylabel={density},
  xmin=-4,xmax=4,ymin=0,ymax=0.62,
  legend pos=north east, legend style={font=\scriptsize, draw=black!30}
  ]
  \addplot[black!45,thick,const plot,domain=-4:4,samples=9]
    coordinates {(-3.5,0.05)(-2.5,0.09)(-1.5,0.22)(-0.5,0.4)(0.5,0.4)(1.5,0.22)(2.5,0.09)(3.5,0.05)};
  \addlegendentry{histogram of $z_m$, moderate $m$}
  \addplot[figred,very thick,domain=-4:4,samples=150] {1/sqrt(2*pi)*exp(-x^2/2)};
  \addlegendentry{$N(0,1)$ limit}
\end{axis}
\end{tikzpicture}
\caption[Central Limit Theorem]{Illustrative histogram of a standardized sample mean for a moderate sample size, compared with the standard normal limiting density (red). The horizontal coordinate is $z=\sqrt{m}(\bar X_m-\mu)/\sigma$.}
\label{fig:ch16-clt}
\end{figure}

\section{Stochastic Differential Equations}
\label{sec:stochastic-differential-equations}

The ordinary differential equations studied in the preceding chapter describe deterministic systems: under existence and uniqueness hypotheses, an initial condition determines a trajectory on its interval of existence. Many biophysical systems, however, experience intrinsic fluctuations, such as thermal noise, fluctuations in molecular counts in a genetic circuit, or synaptic noise in a neural network. To model these effects, we extend $\dot{\bv}=f(\bv)$ by introducing a random driving term. The probability concepts and Gaussian limits developed in this chapter provide the background for this extension.

\subsection{Brownian Motion}

\begin{definition}[Wiener Process]
\label{def:wiener-process}
A \emph{standard Wiener process}, or \emph{standard Brownian motion}, is a real-valued stochastic process $W(t)$, $t\ge0$, with the following properties:
\begin{enumerate}
\item \textit{Initial value}: $W(0)=0$ almost surely.
\item \textit{Independent increments}: it has independent increments: for any $0\leq t_0<t_1<\cdots<t_m$, the random variables $W(t_j)-W(t_{j-1})$, $j=1,\dots,m$, are jointly independent.
\item \textit{Gaussian increments}: for $0\leq s<t$, the increment $W(t)-W(s)$ has distribution $\mathcal{N}(0,t-s)$.
\item \textit{Continuity}: its sample paths $t\mapsto W(t)$ are almost surely continuous.
\end{enumerate}
\end{definition}

\begin{remark}[Continuity without Differentiability]
Brownian paths are almost surely nowhere differentiable, a result we do not prove here. The scale of an increment gives some intuition: over an interval of length $\Delta t$, its standard deviation is $(\Delta t)^{1/2}$, so a typical difference quotient has magnitude of order
\[
\frac{\sqrt{\Delta t}}{\Delta t}=\frac{1}{\sqrt{\Delta t}},
\]
which grows without bound as $\Delta t\to0$. This scaling is a heuristic, not a proof of nowhere differentiability. It explains why $\dot W$ cannot be treated as an ordinary function and why an equation written informally as $\dot{\bv}=f(\bv)+\dot{\bW}$ is not a classical ODE. A rigorous treatment requires stochastic integration.
\end{remark}

Figure~\ref{fig:atlas-brownian-paths} compares independent discretized realizations of Brownian motion.

\begin{figure}[!htbp]
\centering
\begin{tikzpicture}[>=Stealth]
\begin{axis}[
  width=12.5cm,height=6.2cm,
  axis lines=left,
  ticklabel style={font=\small},
  label style={font=\small},
  xlabel={$t$},ylabel={$W_t$},
  xmin=0,xmax=1.03,ymin=-1.4,ymax=2.2,
  xtick={0,0.25,0.5,0.75,1},
  grid=major,grid style={black!10}
]
\addplot[figblue,thin] coordinates {(0.000000,0.000000) (0.003906,0.071894) (0.007812,0.099667) (0.011719,0.161176) (0.015625,0.276561) (0.019531,0.303321) (0.023438,0.274530) (0.027344,0.308431) (0.031250,0.357354) (0.035156,0.294141) (0.039062,0.318164) (0.042969,0.348022) (0.046875,0.360721) (0.050781,0.378277) (0.054688,0.370466) (0.058594,0.371051) (0.062500,0.385740) (0.066406,0.327751) (0.070312,0.381763) (0.074219,0.372351) (0.078125,0.331046) (0.082031,0.286042) (0.085938,0.360965) (0.089844,0.381335) (0.093750,0.303530) (0.097656,0.275398) (0.101562,0.209094) (0.105469,0.158379) (0.109375,0.180587) (0.113281,0.226890) (0.117188,0.279170) (0.121094,0.268315) (0.125000,0.255850) (0.128906,0.201120) (0.132812,0.233901) (0.136719,0.157671) (0.140625,0.161345) (0.144531,0.182987) (0.148438,0.094929) (0.152344,0.174125) (0.156250,0.086354) (0.160156,0.102647) (0.164062,0.134677) (0.167969,0.138649) (0.171875,0.130281) (0.175781,0.164821) (0.179688,0.086271) (0.183594,0.114667) (0.187500,0.111874) (0.191406,0.138953) (0.195312,0.263080) (0.199219,0.174867) (0.203125,0.063584) (0.207031,0.158739) (0.210938,0.105439) (0.214844,0.201024) (0.218750,0.122955) (0.222656,0.242608) (0.226562,0.425723) (0.230469,0.408087) (0.234375,0.487504) (0.238281,0.499184) (0.242188,0.464558) (0.246094,0.438329) (0.250000,0.380200) (0.253906,0.392881) (0.257812,0.340354) (0.261719,0.409444) (0.265625,0.420920) (0.269531,0.352419) (0.273438,0.292158) (0.277344,0.358440) (0.281250,0.412231) (0.285156,0.439973) (0.289062,0.475860) (0.292969,0.483186) (0.296875,0.583534) (0.300781,0.560372) (0.304688,0.573706) (0.308594,0.669594) (0.312500,0.560663) (0.316406,0.596844) (0.320312,0.564603) (0.324219,0.561020) (0.328125,0.548369) (0.332031,0.475480) (0.335938,0.498073) (0.339844,0.475058) (0.343750,0.397360) (0.347656,0.301083) (0.351562,0.328202) (0.355469,0.415408) (0.359375,0.391643) (0.363281,0.320465) (0.367188,0.429462) (0.371094,0.427974) (0.375000,0.430546) (0.378906,0.425323) (0.382812,0.415855) (0.386719,0.460336) (0.390625,0.312285) (0.394531,0.224310) (0.398438,0.206260) (0.402344,0.267705) (0.406250,0.274579) (0.410156,0.080395) (0.414062,0.118034) (0.417969,-0.060501) (0.421875,-0.165487) (0.425781,-0.184497) (0.429688,-0.250202) (0.433594,-0.270557) (0.437500,-0.355921) (0.441406,-0.344766) (0.445312,-0.324919) (0.449219,-0.394724) (0.453125,-0.380942) (0.457031,-0.383763) (0.460938,-0.469585) (0.464844,-0.544962) (0.468750,-0.531145) (0.472656,-0.482526) (0.476562,-0.457292) (0.480469,-0.369933) (0.484375,-0.267920) (0.488281,-0.277632) (0.492188,-0.265407) (0.496094,-0.215602) (0.500000,-0.237320) (0.503906,-0.175597) (0.507812,-0.204590) (0.511719,-0.197919) (0.515625,-0.181113) (0.519531,-0.170213) (0.523438,-0.179235) (0.527344,-0.161139) (0.531250,-0.126996) (0.535156,-0.021619) (0.539062,-0.088090) (0.542969,-0.140985) (0.546875,-0.182087) (0.550781,-0.176136) (0.554688,-0.188205) (0.558594,-0.081948) (0.562500,-0.227149) (0.566406,-0.233622) (0.570312,-0.164457) (0.574219,-0.093145) (0.578125,-0.221449) (0.582031,-0.229593) (0.585938,-0.293614) (0.589844,-0.381834) (0.593750,-0.425104) (0.597656,-0.473408) (0.601562,-0.495901) (0.605469,-0.378299) (0.609375,-0.369934) (0.613281,-0.389069) (0.617188,-0.429705) (0.621094,-0.421121) (0.625000,-0.449874) (0.628906,-0.522024) (0.632812,-0.535298) (0.636719,-0.635955) (0.640625,-0.572943) (0.644531,-0.687222) (0.648438,-0.587140) (0.652344,-0.451273) (0.656250,-0.443199) (0.660156,-0.377786) (0.664062,-0.307570) (0.667969,-0.347673) (0.671875,-0.348133) (0.675781,-0.411539) (0.679688,-0.490655) (0.683594,-0.511667) (0.687500,-0.476269) (0.691406,-0.493906) (0.695312,-0.523734) (0.699219,-0.571607) (0.703125,-0.628278) (0.707031,-0.516322) (0.710938,-0.513697) (0.714844,-0.521439) (0.718750,-0.585341) (0.722656,-0.568698) (0.726562,-0.556772) (0.730469,-0.572136) (0.734375,-0.494649) (0.738281,-0.466995) (0.742188,-0.517893) (0.746094,-0.593066) (0.750000,-0.598301) (0.753906,-0.552629) (0.757812,-0.579400) (0.761719,-0.509875) (0.765625,-0.553263) (0.769531,-0.612072) (0.773438,-0.579023) (0.777344,-0.469131) (0.781250,-0.525759) (0.785156,-0.489410) (0.789062,-0.438775) (0.792969,-0.467512) (0.796875,-0.479576) (0.800781,-0.553124) (0.804688,-0.629102) (0.808594,-0.543822) (0.812500,-0.531688) (0.816406,-0.504896) (0.820312,-0.538428) (0.824219,-0.446751) (0.828125,-0.431713) (0.832031,-0.554976) (0.835938,-0.482365) (0.839844,-0.503309) (0.843750,-0.530958) (0.847656,-0.628338) (0.851562,-0.544194) (0.855469,-0.557405) (0.859375,-0.615450) (0.863281,-0.507438) (0.867188,-0.503442) (0.871094,-0.575805) (0.875000,-0.616645) (0.878906,-0.635226) (0.882812,-0.694043) (0.886719,-0.680323) (0.890625,-0.749864) (0.894531,-0.707605) (0.898438,-0.778608) (0.902344,-0.700066) (0.906250,-0.804586) (0.910156,-0.893449) (0.914062,-0.811714) (0.917969,-0.735168) (0.921875,-0.809518) (0.925781,-0.833310) (0.929688,-0.856793) (0.933594,-0.864099) (0.937500,-0.908484) (0.941406,-0.889948) (0.945312,-0.905708) (0.949219,-0.916522) (0.953125,-0.974218) (0.957031,-1.024776) (0.960938,-1.200692) (0.964844,-1.166844) (0.968750,-1.107111) (0.972656,-1.082881) (0.976562,-1.115680) (0.980469,-1.058859) (0.984375,-1.035515) (0.988281,-1.119239) (0.992188,-1.118584) (0.996094,-1.126853) (1.000000,-1.262117)};
\addplot[figred,thin] coordinates {(0.000000,0.000000) (0.003906,-0.029506) (0.007812,-0.065051) (0.011719,-0.028742) (0.015625,-0.096435) (0.019531,-0.150042) (0.023438,-0.154057) (0.027344,-0.168918) (0.031250,-0.175222) (0.035156,-0.111769) (0.039062,-0.149804) (0.042969,-0.256598) (0.046875,-0.215331) (0.050781,-0.184882) (0.054688,-0.136478) (0.058594,-0.206721) (0.062500,-0.323506) (0.066406,-0.295969) (0.070312,-0.286195) (0.074219,-0.287495) (0.078125,-0.240818) (0.082031,-0.220457) (0.085938,-0.252248) (0.089844,-0.299812) (0.093750,-0.394050) (0.097656,-0.343627) (0.101562,-0.357341) (0.105469,-0.445727) (0.109375,-0.444332) (0.113281,-0.461245) (0.117188,-0.417181) (0.121094,-0.500581) (0.125000,-0.463202) (0.128906,-0.439790) (0.132812,-0.430240) (0.136719,-0.400264) (0.140625,-0.431789) (0.144531,-0.388948) (0.148438,-0.337485) (0.152344,-0.372051) (0.156250,-0.461285) (0.160156,-0.440345) (0.164062,-0.441948) (0.167969,-0.385368) (0.171875,-0.409401) (0.175781,-0.456165) (0.179688,-0.458066) (0.183594,-0.504559) (0.187500,-0.511522) (0.191406,-0.516361) (0.195312,-0.541767) (0.199219,-0.561876) (0.203125,-0.598022) (0.207031,-0.619152) (0.210938,-0.626959) (0.214844,-0.671837) (0.218750,-0.730047) (0.222656,-0.724819) (0.226562,-0.642867) (0.230469,-0.582404) (0.234375,-0.552744) (0.238281,-0.603631) (0.242188,-0.672372) (0.246094,-0.637026) (0.250000,-0.578936) (0.253906,-0.464766) (0.257812,-0.381258) (0.261719,-0.292269) (0.265625,-0.324989) (0.269531,-0.314480) (0.273438,-0.264114) (0.277344,-0.264684) (0.281250,-0.248511) (0.285156,-0.195882) (0.289062,-0.061264) (0.292969,-0.043557) (0.296875,0.034533) (0.300781,0.063224) (0.304688,0.048614) (0.308594,-0.004877) (0.312500,0.037370) (0.316406,0.119303) (0.320312,0.105191) (0.324219,0.061354) (0.328125,0.134040) (0.332031,0.204421) (0.335938,0.215187) (0.339844,0.178180) (0.343750,0.142985) (0.347656,0.083378) (0.351562,0.060888) (0.355469,0.164955) (0.359375,0.175233) (0.363281,0.143365) (0.367188,0.253995) (0.371094,0.250465) (0.375000,0.214465) (0.378906,0.201689) (0.382812,0.186901) (0.386719,0.224599) (0.390625,0.167184) (0.394531,0.107884) (0.398438,0.027083) (0.402344,0.138455) (0.406250,0.201816) (0.410156,0.175921) (0.414062,0.225826) (0.417969,0.086679) (0.421875,0.026409) (0.425781,-0.010201) (0.429688,0.067482) (0.433594,0.054060) (0.437500,0.114557) (0.441406,0.115057) (0.445312,0.200676) (0.449219,0.154276) (0.453125,0.141920) (0.457031,0.135491) (0.460938,0.233644) (0.464844,0.149857) (0.468750,0.111019) (0.472656,0.079022) (0.476562,0.136081) (0.480469,0.138743) (0.484375,0.258287) (0.488281,0.235928) (0.492188,0.275948) (0.496094,0.211159) (0.500000,0.214171) (0.503906,0.183368) (0.507812,0.210278) (0.511719,0.147170) (0.515625,0.182961) (0.519531,0.171410) (0.523438,0.112894) (0.527344,0.099114) (0.531250,0.214013) (0.535156,0.335778) (0.539062,0.305042) (0.542969,0.347105) (0.546875,0.334276) (0.550781,0.428011) (0.554688,0.299347) (0.558594,0.327492) (0.562500,0.341076) (0.566406,0.355604) (0.570312,0.408944) (0.574219,0.460406) (0.578125,0.549397) (0.582031,0.583684) (0.585938,0.655660) (0.589844,0.656731) (0.593750,0.696730) (0.597656,0.650395) (0.601562,0.732357) (0.605469,0.769263) (0.609375,0.792091) (0.613281,0.787495) (0.617188,0.773844) (0.621094,0.784915) (0.625000,0.887730) (0.628906,0.881640) (0.632812,0.867941) (0.636719,0.877520) (0.640625,0.968315) (0.644531,0.894307) (0.648438,0.785656) (0.652344,0.900927) (0.656250,0.921730) (0.660156,1.029922) (0.664062,1.089172) (0.667969,1.062291) (0.671875,1.151485) (0.675781,1.067422) (0.679688,1.137168) (0.683594,1.140709) (0.687500,1.175807) (0.691406,1.223792) (0.695312,1.171859) (0.699219,1.272209) (0.703125,1.208656) (0.707031,1.209317) (0.710938,1.298264) (0.714844,1.344554) (0.718750,1.249169) (0.722656,1.303576) (0.726562,1.314609) (0.730469,1.427168) (0.734375,1.499224) (0.738281,1.584170) (0.742188,1.588172) (0.746094,1.744125) (0.750000,1.755908) (0.753906,1.652448) (0.757812,1.640733) (0.761719,1.723463) (0.765625,1.781676) (0.769531,1.771427) (0.773438,1.682514) (0.777344,1.575211) (0.781250,1.517077) (0.785156,1.559594) (0.789062,1.452903) (0.792969,1.511255) (0.796875,1.529966) (0.800781,1.509829) (0.804688,1.380237) (0.808594,1.473233) (0.812500,1.586517) (0.816406,1.617732) (0.820312,1.682910) (0.824219,1.745530) (0.828125,1.724981) (0.832031,1.681721) (0.835938,1.727572) (0.839844,1.729809) (0.843750,1.684408) (0.847656,1.813063) (0.851562,1.883828) (0.855469,1.899424) (0.859375,1.907413) (0.863281,1.908964) (0.867188,1.894269) (0.871094,1.894773) (0.875000,1.891761) (0.878906,1.846694) (0.882812,1.835272) (0.886719,1.862396) (0.890625,1.847036) (0.894531,1.864544) (0.898438,1.938020) (0.902344,2.006924) (0.906250,1.944791) (0.910156,1.973791) (0.914062,1.919343) (0.917969,1.903685) (0.921875,1.882549) (0.925781,1.800078) (0.929688,1.880092) (0.933594,1.781018) (0.937500,1.761642) (0.941406,1.770016) (0.945312,1.776495) (0.949219,1.791739) (0.953125,1.819722) (0.957031,1.777823) (0.960938,1.835601) (0.964844,1.903809) (0.968750,1.960148) (0.972656,1.932963) (0.976562,1.972191) (0.980469,2.030314) (0.984375,1.886979) (0.988281,1.835040) (0.992188,1.912010) (0.996094,1.848262) (1.000000,1.880328)};
\addplot[figgreen,thin] coordinates {(0.000000,0.000000) (0.003906,-0.025111) (0.007812,-0.005208) (0.011719,0.057699) (0.015625,0.103643) (0.019531,0.175262) (0.023438,0.300538) (0.027344,0.364271) (0.031250,0.297284) (0.035156,0.293449) (0.039062,0.276558) (0.042969,0.297778) (0.046875,0.352644) (0.050781,0.336403) (0.054688,0.341060) (0.058594,0.390238) (0.062500,0.399815) (0.066406,0.477564) (0.070312,0.505813) (0.074219,0.594553) (0.078125,0.651313) (0.082031,0.720662) (0.085938,0.718457) (0.089844,0.778104) (0.093750,0.712839) (0.097656,0.785709) (0.101562,0.774404) (0.105469,0.772335) (0.109375,0.736525) (0.113281,0.756319) (0.117188,0.810040) (0.121094,0.774600) (0.125000,0.772764) (0.128906,0.834047) (0.132812,0.813085) (0.136719,0.699305) (0.140625,0.766272) (0.144531,0.770126) (0.148438,0.771713) (0.152344,0.898191) (0.156250,0.879438) (0.160156,0.888041) (0.164062,0.973387) (0.167969,1.106829) (0.171875,1.238212) (0.175781,1.237260) (0.179688,1.187050) (0.183594,1.240801) (0.187500,1.131164) (0.191406,1.132831) (0.195312,1.152186) (0.199219,1.152715) (0.203125,1.123587) (0.207031,1.175339) (0.210938,1.197038) (0.214844,1.120756) (0.218750,1.164631) (0.222656,1.132019) (0.226562,1.170531) (0.230469,1.247165) (0.234375,1.301283) (0.238281,1.277580) (0.242188,1.397103) (0.246094,1.405889) (0.250000,1.368024) (0.253906,1.288213) (0.257812,1.336854) (0.261719,1.276878) (0.265625,1.348805) (0.269531,1.389297) (0.273438,1.331032) (0.277344,1.335090) (0.281250,1.400331) (0.285156,1.282302) (0.289062,1.294682) (0.292969,1.122304) (0.296875,1.109572) (0.300781,1.082782) (0.304688,1.004933) (0.308594,0.931148) (0.312500,0.881930) (0.316406,0.901384) (0.320312,0.930065) (0.324219,0.970049) (0.328125,1.015457) (0.332031,1.052842) (0.335938,1.016752) (0.339844,1.014984) (0.343750,0.988403) (0.347656,1.044343) (0.351562,0.925725) (0.355469,0.878376) (0.359375,0.810663) (0.363281,0.758180) (0.367188,0.863064) (0.371094,0.797596) (0.375000,0.932530) (0.378906,0.903269) (0.382812,0.869623) (0.386719,0.853793) (0.390625,0.861174) (0.394531,0.844708) (0.398438,0.838100) (0.402344,0.881782) (0.406250,1.028804) (0.410156,0.960426) (0.414062,0.926870) (0.417969,0.902357) (0.421875,0.966700) (0.425781,0.947902) (0.429688,0.978309) (0.433594,1.044760) (0.437500,1.044206) (0.441406,1.006303) (0.445312,1.018223) (0.449219,1.023594) (0.453125,1.003573) (0.457031,0.978695) (0.460938,0.923673) (0.464844,0.843002) (0.468750,0.818019) (0.472656,0.832027) (0.476562,0.812912) (0.480469,0.749909) (0.484375,0.755257) (0.488281,0.773499) (0.492188,0.821470) (0.496094,0.822808) (0.500000,0.804076) (0.503906,0.812094) (0.507812,0.876290) (0.511719,0.736135) (0.515625,0.825891) (0.519531,0.728027) (0.523438,0.644779) (0.527344,0.716917) (0.531250,0.686794) (0.535156,0.736577) (0.539062,0.569293) (0.542969,0.488366) (0.546875,0.465007) (0.550781,0.482931) (0.554688,0.447180) (0.558594,0.470264) (0.562500,0.486013) (0.566406,0.484168) (0.570312,0.517742) (0.574219,0.412646) (0.578125,0.392537) (0.582031,0.421263) (0.585938,0.424241) (0.589844,0.474981) (0.593750,0.500363) (0.597656,0.461650) (0.601562,0.452201) (0.605469,0.482847) (0.609375,0.472603) (0.613281,0.520173) (0.617188,0.577951) (0.621094,0.568870) (0.625000,0.569530) (0.628906,0.656062) (0.632812,0.636640) (0.636719,0.672636) (0.640625,0.626296) (0.644531,0.590166) (0.648438,0.640416) (0.652344,0.636585) (0.656250,0.601023) (0.660156,0.497326) (0.664062,0.456873) (0.667969,0.351289) (0.671875,0.321702) (0.675781,0.345032) (0.679688,0.293851) (0.683594,0.398429) (0.687500,0.408629) (0.691406,0.374026) (0.695312,0.417987) (0.699219,0.427224) (0.703125,0.517577) (0.707031,0.581326) (0.710938,0.665200) (0.714844,0.642112) (0.718750,0.542446) (0.722656,0.504294) (0.726562,0.513517) (0.730469,0.572497) (0.734375,0.589831) (0.738281,0.553430) (0.742188,0.570038) (0.746094,0.610921) (0.750000,0.732581) (0.753906,0.773528) (0.757812,0.837833) (0.761719,0.895476) (0.765625,0.852360) (0.769531,0.870582) (0.773438,0.892371) (0.777344,0.767355) (0.781250,0.723896) (0.785156,0.738528) (0.789062,0.827874) (0.792969,0.786093) (0.796875,0.739173) (0.800781,0.746967) (0.804688,0.863414) (0.808594,0.847187) (0.812500,0.827386) (0.816406,0.762280) (0.820312,0.815711) (0.824219,0.797078) (0.828125,0.861175) (0.832031,0.745486) (0.835938,0.736451) (0.839844,0.889338) (0.843750,0.934403) (0.847656,0.953715) (0.851562,0.956925) (0.855469,0.852220) (0.859375,0.793411) (0.863281,0.775036) (0.867188,0.796473) (0.871094,0.790292) (0.875000,0.808401) (0.878906,0.822130) (0.882812,0.789090) (0.886719,0.668311) (0.890625,0.682052) (0.894531,0.677788) (0.898438,0.667126) (0.902344,0.701464) (0.906250,0.699755) (0.910156,0.627214) (0.914062,0.559267) (0.917969,0.536071) (0.921875,0.591080) (0.925781,0.648641) (0.929688,0.601301) (0.933594,0.613979) (0.937500,0.549664) (0.941406,0.529575) (0.945312,0.500962) (0.949219,0.462865) (0.953125,0.486680) (0.957031,0.639620) (0.960938,0.688515) (0.964844,0.570401) (0.968750,0.670939) (0.972656,0.624456) (0.976562,0.649979) (0.980469,0.641179) (0.984375,0.665896) (0.988281,0.726512) (0.992188,0.772373) (0.996094,0.763951) (1.000000,0.743884)};
\end{axis}
\end{tikzpicture}
\caption{Three reproducibly simulated paths start at zero and use independent Gaussian increments of variance $\Delta t=1/256$. Straight segments interpolate the sampled values for display. These finite-resolution drawings do not imply differentiability of Brownian paths and are not experimental observations.}
\label{fig:atlas-brownian-paths}
\end{figure}

\subsection{The Langevin Equation}

\begin{definition}[It\^o Stochastic Differential Equation]
\label{def:langevin-equation}
An It\^o \emph{Langevin equation} has the form
\begin{equation}
\dd\bv(t)=f(\bv(t))\,\dd t+\sigma(\bv(t))\,\dd\bW(t),
\end{equation}
where $\bv(t)\in\R^d$, $f:\R^d\to\R^d$ is the deterministic \emph{drift}, $\bW$ is an $m$-dimensional Wiener process with independent components, and $\sigma:\R^d\to\R^{d\times m}$ is the noise-amplitude matrix. The diffusion covariance is $\sigma\sigma^{\top}$. The equation is interpreted in integral form:
\begin{equation}
\bv(t)=\bv(0)+\int_0^t f(\bv(s))\,\dd s
+\int_0^t\sigma(\bv(s))\,\dd\bW(s).
\end{equation}

The last term is an \emph{It\^o stochastic integral}. For a suitable adapted integrand, it is constructed as a limit of non-anticipating sums; for regular integrands, these take the left-endpoint form
\[
\sum_j\sigma(\bv(t_j))\big(\bW(t_{j+1})-\bW(t_j)\big).
\]

Unlike ordinary Riemann integration, the evaluation convention matters. A different convention, such as Stratonovich integration, generally changes the drift when $\sigma$ depends on the state.
\end{definition}

For example, global Lipschitz and linear-growth assumptions on $f$ and $\sigma$, together with a square-integrable initial state independent of future Brownian increments, ensure a unique global strong solution. We do not develop the existence theory here; these assumptions indicate the role played by regularity conditions, analogous to those required for deterministic ODEs.

\begin{example}[The Ornstein--Uhlenbeck Process]
\label{ex:ornstein-uhlenbeck}
In one dimension, linear relaxation toward zero with rate $\gamma>0$ and constant noise amplitude $\sigma>0$ gives
\begin{equation}
\dd v=-\gamma v\,\dd t+\sigma\,\dd W.
\end{equation}

This is the \emph{Ornstein--Uhlenbeck process}, the classical Langevin model for a particle velocity subject to viscous friction and random molecular impacts. Its drift is the same linear relaxation encountered in deterministic first-order equations, while its noise term represents the fluctuating forcing. It also describes the subthreshold dynamics of a leaky integrate-and-fire neuron driven by Gaussian noise after shifting the deterministic equilibrium to zero; a full spiking model additionally requires a threshold and reset rule.
\end{example}

\section{Diffusion Equations and Stationary Probability Flows}

\subsection{The Fokker--Planck Equation and Infinitesimal Moments}


Instead of following a single random trajectory, we may study the time-dependent probability density $p(v,t)$ of the state across all noise realizations. Under suitable assumptions, this density satisfies a deterministic PDE.

\begin{theorem}[Fokker--Planck Equation for an It\^o Diffusion]
\label{thm:fokker-planck}
Let $v(t)$ solve the one-dimensional It\^o equation
\[
\dd v=f(v)\,\dd t+\sigma(v)\,\dd W.
\]

Assume that the coefficients and the probability density $p(v,t)$ are sufficiently regular for the following derivatives to exist classically. Then
\begin{equation}
\frac{\partial p}{\partial t}
=-\frac{\partial}{\partial v}\big[f(v)p\big]
+\frac12\frac{\partial^2}{\partial v^2}\big[\sigma(v)^2p\big].
\end{equation}

With less regularity, the equation is interpreted weakly; an arbitrary stochastic differential equation need not have a smooth density.
\end{theorem}

\begin{proof}[Heuristic Derivation via the Kramers--Moyal Expansion]
\leavevmode\par\nobreak
\textit{Step 1: express evolution through transition moments.} For a Markov process, the conditional distribution of a future state depends on the past only through the present state. Under conditions that justify the moment expansion, the forward equation can formally be written as

\begin{equation}
\frac{\partial p}{\partial t}
=\sum_{n=1}^{\infty}\frac{(-1)^n}{n!}
\frac{\partial^n}{\partial v^n}\big[D_n(v)p\big].
\end{equation}
The coefficients in this expansion are defined by
\begin{equation}
D_n(v):=\lim_{\Delta t\downarrow0}\frac{1}{\Delta t}
\E\big[(v(t+\Delta t)-v(t))^n\mid v(t)=v\big].
\end{equation}

These are conditional raw moments of the increment per unit time, not moments centered about its conditional mean. The Kramers--Moyal series is used here as a formal tool, rather than asserted as a convergent representation for every Markov process.

\textit{Step 2: compute the infinitesimal moments.} Over a short interval, freezing the coefficients at the current state gives the leading approximation

\begin{align}
v(t+\Delta t)-v(t)
&=\int_t^{t+\Delta t}f(v(s))\,\dd s
+\int_t^{t+\Delta t}\sigma(v(s))\,\dd W(s)\notag\\
&\approx f(v)\Delta t+\sigma(v)\Delta W,
\end{align}
where $\Delta W\sim\mathcal{N}(0,\Delta t)$ independently of the current state. All expectations below are conditional on $v(t)=v$. Since $\E[\Delta W]=0$,
\begin{equation}
\begin{aligned}
D_1(v) &= \lim_{\Delta t\downarrow0}\frac{f(v)\Delta t}{\Delta t} \\
&= f(v).
\end{aligned}
\end{equation}
For the second moment, the frozen-coefficient approximation gives
\begin{align}
\E\big[(v(t+\Delta t)-v(t))^2\mid v(t)=v\big]
&\approx f(v)^2\Delta t^2
+2f(v)\sigma(v)\Delta t\,\E[\Delta W]\notag\\
&\quad+\sigma(v)^2\E[(\Delta W)^2]\notag\\
&=f(v)^2\Delta t^2+\sigma(v)^2\Delta t,
\end{align}
because
\[
\begin{aligned}
\E[(\Delta W)^2] &= \var(\Delta W) \\
&= \Delta t.
\end{aligned}
\]
Dividing by $\Delta t$ and taking the limit yields
\begin{equation}
D_2(v)=\sigma(v)^2.
\end{equation}

For $n\ge3$, the higher increment moments are $o(\Delta t)$ under suitable moment bounds, so $D_n(v)=0$. The Gaussian scaling underlying this fact is
\[
\E[|\Delta W|^n]=O((\Delta t)^{n/2});
\]
for example,
\[
\begin{gathered}
\E[(\Delta W)^3]=0 \\
\E[(\Delta W)^4]=3\Delta t^2.
\end{gathered}
\]

The neglected terms must be controlled at the level of conditional moments to turn this heuristic into a proof.

\textit{Step 3: retain the first two terms.} Substitution into the formal expansion gives

\begin{align}
\frac{\partial p}{\partial t}
&=-\frac{\partial}{\partial v}(D_1p)
+\frac12\frac{\partial^2}{\partial v^2}(D_2p)\notag\\
&=-\frac{\partial}{\partial v}\big[f(v)p\big]
+\frac12\frac{\partial^2}{\partial v^2}\big[\sigma(v)^2p\big],
\end{align}
which is the Fokker--Planck equation. A rigorous weak derivation uses It\^o's formula on smooth compactly supported test functions in place of the formal series.
\end{proof}

\subsection{Probability Currents and Stationary Densities}

\begin{remark}[A Continuity Equation with Diffusive Current]
The Fokker--Planck equation is a continuity equation of the type derived in Example~\ref{ex:continuity-equation}:
\begin{equation}
\frac{\partial p}{\partial t}+\frac{\partial J}{\partial v}=0,
\qquad J(v,t)=f(v)p(v,t)-\frac12\frac{\partial}{\partial v}\big[\sigma(v)^2p(v,t)\big].
\end{equation}

The first term in the probability current is deterministic transport by the drift; the second is a diffusive flux. For constant $\sigma$ and zero drift, this reduces to the heat equation with $D=\sigma^2/2$, studied in Section~\ref{sec:pde-methods}. Thus random trajectories produce the same diffusion equation for a probability density that describes heat conduction or molecular spreading in other contexts. Conservation of total probability requires suitable boundary conditions, such as zero flux at finite boundaries or vanishing current at infinity.
\end{remark}

\begin{example}[Noise in a Genetic Circuit or a Neuron]
For the Ornstein--Uhlenbeck process of Example~\ref{ex:ornstein-uhlenbeck}, the density satisfies
\begin{equation}
\frac{\partial p}{\partial t}
=\gamma\frac{\partial}{\partial v}(vp)
+\frac{\sigma^2}{2}\frac{\partial^2p}{\partial v^2}.
\end{equation}
On the whole real line, the normalized stationary density is
\begin{equation}
p_\infty(v)=\sqrt{\frac{\gamma}{\pi\sigma^2}}
\exp\left(-\frac{\gamma v^2}{\sigma^2}\right),
\qquad \var(v)=\frac{\sigma^2}{2\gamma}.
\end{equation}
Indeed,
\[
p_\infty'=-(2\gamma v/\sigma^2)p_\infty,
\]
so the stationary current
\[
J_\infty=-\gamma v p_\infty-(\sigma^2/2)p_\infty'
\]
vanishes identically. The restoring drift narrows the distribution, while noise broadens it. This balance models fluctuations of a subthreshold membrane potential around its equilibrium, or small fluctuations of gene expression around a stable mean. In the latter interpretation, $v$ is a deviation from the mean, not an unrestricted physical molecular count; threshold/reset dynamics or positivity constraints require additional modeling.
\end{example}